% ECONOMICS_PROBLEM: {"id": "I38-537", "title": "Preventing persistent homelessness", "jel_code": "I38", "source_id": "OP-0469"}
% !TeX program = xelatex
\documentclass[11pt,a4paper]{article}
\usepackage[margin=23mm]{geometry}
\usepackage{fontspec}
\setmainfont{Latin Modern Roman}
\usepackage{amsmath,amssymb,mathtools}
\usepackage{xcolor,enumitem,booktabs,longtable,array,xurl,hyperref}
\definecolor{muted}{HTML}{53636C}
\hypersetup{colorlinks=true,urlcolor=blue,linkcolor=blue}
\setlength{\parindent}{0pt}
\setlength{\parskip}{5pt}
\newcommand{\R}{\mathbb R}
\newcommand{\E}{\mathbb E}
\newcommand{\N}{\mathbb N}
\newcommand{\ind}{\mathbf 1}
\newcommand{\Var}{\operatorname{Var}}
\newcommand{\Cov}{\operatorname{Cov}}
\newcommand{\supp}{\operatorname{supp}}
\DeclareMathOperator*{\argmin}{arg\,min}
\DeclareMathOperator*{\argmax}{arg\,max}
\newcommand{\entrymeta}[3]{{\sffamily\small\color{muted}\textbf{\texttt{#1}}\quad|\quad #2\par #3\par}}
\newcommand{\Problemref}[1]{\texttt{#1}}
\newcommand{\JELref}[2]{\texttt{#2} (JEL \texttt{#1})}
\begin{document}
\section*{Preventing persistent homelessness}
\textbf{Problem ID:} \texttt{I38-537}\quad \textbf{JEL:} \texttt{I38}\par
% BEGIN ECONOMICS STATEMENT
\label{problem:OP-0469}
{\sffamily\small\color{muted}Original subject group: Urban, housing and spatial economics\par}
\label{definitions:problem:30007751}
\textbf{Audit classification:} Published research agenda. Specified feasible measurable policy classes, capacities and objective units; known-response optimization is routine.
\subsection*{Definitions and precise research target}
\textit{Editorial formalization.} Consider adults entering a coordinated homelessness-assessment system in one city during a specified year. For person \(i\), let baseline covariates \(X_i\) include housing history, health, income, and family composition. A feasible allocation rule \(d(X_i)\in\{0,1,2\}\) assigns existing services, temporary rental assistance, or permanent supportive housing. Let \(H_i(d)\) be days without stable housing during the next twenty-four months, \(Q_i(d)\) health-related utility, and \(C_i(d)\) real service cost. For a per-eligible-person budget \(B\) and health weight \(\lambda>0\), define
\[d^*\in\arg\max_d E[\lambda Q_i(d(X_i))-H_i(d(X_i))]\quad\text{subject to}\quad E[C_i(d(X_i))]\leq B.\]
Restrict rules to a declared feasible class with finite outcome means. Establish attainment or report the supremum and near-optimal rules. State capacity and equity constraints before optimization. Identify treatment-response heterogeneity through randomized offers or a justified allocation discontinuity, distinguishing predicted homelessness risk from predicted treatment benefit. Follow people across shelters, municipalities, and stable housing exits, with explicit missing-data assumptions. Include crowd-out and housing-market spillovers where program scale makes them relevant.

Let covariates \(X\) lie in a declared measurable space and policies \(d\) in a nonempty finite class \(\mathcal D\) of measurable maps to \(\{0,1,2\}\), or use an explicit randomized class \(d_a(X)\geq0\), \(\sum_ad_a(X)=1\). \(H_i(a)\) counts days outside an operational stable-housing definition during twenty-four months, \(Q_i(a)\) is quality-adjusted health over that window, and \(C_i(a)\) incremental real service cost. \(\lambda\) has units housing-days per quality-adjusted year; the objective is a stated composite rather than unnormalized utility. Include a feasible baseline whose cost does not exceed \(B\), finite means, capacities and equity restrictions. A finite feasible policy class attains an optimum. Identify conditional benefits \(E[\lambda Q(a)-H(a)\mid X]\), not only baseline risk, while addressing offers versus receipt and crowd-out. All expectations use a stated baseline population and probability law. Identification means every admissible data-generating model with the same observed-data law yields the same displayed target; otherwise report the identified set. Exact equations, horizons and assumptions are editorial, rather than source-given canonical conjectures.
\subsection*{Variants and assumption changes}
\begin{enumerate}
\item \textbf{Editorial randomized-rule variant.} Permit assignment probabilities \(d_a(X)\) with \(\sum_a d_a(X)=1\) and \(E[d_a(X)]\leq\kappa_a\). Randomized ties can satisfy budgets/capacities differently.
\item \textbf{Editorial dynamic-service variant.} Reassess every six months under sequential randomized offers; define history-dependent rules, recurrence and cumulative budgets. One-time effects do not identify this policy.
\end{enumerate}
\subsection*{References and source audit}
\textbf{1.} Explicit targeting/matching agenda distinguishes predicted risk from treatment benefit. Locator: J-PAL report, July 2019; targeting and matching discussion, printed pp.45–46. Full text or primary publication page inspected. URL: \url{https://www.povertyactionlab.org/sites/default/files/research-paper/reducing-preventing-homelessness-review-research-agenda.pdf}.
\begin{enumerate}\raggedright
\item\hypertarget{definitions:source:30007751:1}{} William N. Evans; David C. Phillips; Krista Ruffini. 2019. \emph{Reducing and Preventing Homelessness: A Review of the Evidence and Charting a Research Agenda}. J-PAL report, July 2019; targeting and matching discussion, printed pp.45–46.
\url{https://www.povertyactionlab.org/sites/default/files/research-paper/reducing-preventing-homelessness-review-research-agenda.pdf}.
DOI: \url{https://doi.org/10.3386/w26232}.
\end{enumerate}

\subsection*{Common conventions: Source mathematical conventions}

These conventions apply to each entry unless explicitly replaced.

\textbf{Models and identification.} A primitive model \(m\) specifies all probability laws, feasible choices, information, equilibrium and measurement rules used by its target. The admissible class \(\mathcal M\) is fixed independently of outcomes. If \(P_O(m)\) is its observable probability law and \(\tau(m)\) the target, its identified set at observable law \(P\) is
\[
 \mathcal I_\tau(P)=\{\tau(m):m\in\mathcal M,\ P_O(m)=P\}.
\]
Point identification means that this set is a singleton. An empty set signals incompatibility of data and assumptions. Coordinate bounds need not describe the joint set. Bounds are sharp only if every retained target is attainable or arbitrarily approachable by compatible models. Finite bounds require explicit restrictions on unobserved quantities; unrestricted confounding, hidden wealth, extrapolation or arbitrary equilibrium selection can yield uninformative sets.

\textbf{Causal interventions.} A potential outcome \(Y_i(p)\) is the outcome under a complete policy regime \(p\), including timing, financing, information and market spillovers. The target population and units are fixed before treatment. With interference, use the complete assignment vector or a justified exposure mapping; individual no-interference notation alone cannot identify an economy-wide intervention. A difference of mean potential outcomes does not require the cross-world joint distribution. Conditioning on post-treatment migration, survival, enrollment or employment changes the estimand and needs separate justification.

\textbf{Statistical statements.} An estimator, confidence set, forecast or test specifies a sampling law, model class and loss/coverage criterion. Uniform \((1-\alpha)\)-coverage means \(\inf_{m\in\mathcal M}\Pr_m\{\tau(m)\in C_n\}\geq1-\alpha\), or an explicitly asymptotic version. The whole parameter space trivially has coverage, so informativeness requires a stated width, risk or power objective. A forecasting improvement is relative to a named benchmark, information set and evaluation population.

\textbf{Equilibrium and optimization.} An equilibrium comprises feasible plans, optimality, beliefs where needed, budget constraints and all market/resource clearing. The primitives or a stated benchmark define each correspondence \(\mathcal E(p;m)\). Nonemptiness is assumed or proved, never inferred just from compact policy sets. A selected-equilibrium target uses a declared selection rule; a robust target quantifies over all permitted equilibria. Use suprema or infima unless attainment is established. An \(\varepsilon\)-optimal policy is feasible and within \(\varepsilon\) of the finite optimum value. Report an empty feasible set separately.

\textbf{Welfare and accounting.} Normative utility, interpersonal normalization, social weights, numeraire, discounting and the use of public receipts are primitives. Behavioral utility need not coincide with the welfare criterion. Household welfare includes ownership income and rents once; adding profits again to compensating variation generally double-counts them. A monetary equivalent is defined by an explicit transfer and price convention, with monotonicity and range conditions. Average effects, mechanism contributions, transfers and welfare losses are different targets. Infinite sums require convergence and dynamic feasible plans require terminal or no-Ponzi conditions.

\textbf{Source versions.} A working-paper date, revision date and journal publication year are distinguished. A DOI for a journal version is not automatically the DOI of an earlier manuscript. Locators refer to the stated version and printed pages where specified. A metadata-only check does not verify a claim about a full-text conclusion. Every entry's source subsection narrows the provenance claim accordingly.
% END ECONOMICS STATEMENT
\end{document}
